\documentclass[10pt,a4paper]{amsart}
\usepackage{amsmath,amssymb,amsthm,graphicx}
% Retained preamble for the Stacks Project chapter {CH} (shared project preamble).
% The chapter file pinned-original-source.tex only carries \input{{preamble}} on its first line,
% so no preamble text lives inside the chapter itself.  The Stacks Project uses one single shared
% preamble for every chapter: the snapshot stages it three times (sources/stacks-chow/preamble.tex,
% sources/stacks-etale-cohomology/preamble.tex and focused-seeds/AGP/sources/stacks-snapshot-20261011/
% preamble.tex) and those copies are identical in size.  The block below is copied mechanically,
% line for line, from sources/stacks-chow/preamble.tex lines 14-18 and 154-217: the Xy-pic setup and
% the native theorem declarations together with the mathematical macros.  Nothing was added, removed
% or redefined; the wrapper class and the amsmath/amssymb/amsthm/graphicx base packages are supplied
% by the assembler.
% For commutative diagrams we use Xy-pic
% Chapter: stacks-spaces-more-morphisms.tex -- More on Morphisms of Spaces.
\usepackage[all]{xy}

% We use 2cell for 2-commutative diagrams.
\xyoption{2cell}

% Theorem environments.
%
\theoremstyle{plain}
\newtheorem{theorem}[subsection]{Theorem}
\newtheorem{proposition}[subsection]{Proposition}
\newtheorem{lemma}[subsection]{Lemma}

\theoremstyle{definition}
\newtheorem{definition}[subsection]{Definition}
\newtheorem{example}[subsection]{Example}
\newtheorem{exercise}[subsection]{Exercise}
\newtheorem{situation}[subsection]{Situation}

\theoremstyle{remark}
\newtheorem{remark}[subsection]{Remark}
\newtheorem{remarks}[subsection]{Remarks}

\numberwithin{equation}{subsection}

% Macros
%
\def\lim{\mathop{\mathrm{lim}}\nolimits}
\def\colim{\mathop{\mathrm{colim}}\nolimits}
\def\Spec{\mathop{\mathrm{Spec}}}
\def\Hom{\mathop{\mathrm{Hom}}\nolimits}
\def\Ext{\mathop{\mathrm{Ext}}\nolimits}
\def\SheafHom{\mathop{\mathcal{H}\!\mathit{om}}\nolimits}
\def\SheafExt{\mathop{\mathcal{E}\!\mathit{xt}}\nolimits}
\def\Sch{\mathit{Sch}}
\def\Mor{\mathop{\mathrm{Mor}}\nolimits}
\def\Ob{\mathop{\mathrm{Ob}}\nolimits}
\def\Sh{\mathop{\mathit{Sh}}\nolimits}
\def\NL{\mathop{N\!L}\nolimits}
\def\CH{\mathop{\mathrm{CH}}\nolimits}
\def\proetale{{pro\text{-}\acute{e}tale}}
\def\etale{{\acute{e}tale}}
\def\QCoh{\mathit{QCoh}}
\def\Ker{\mathop{\mathrm{Ker}}}
\def\Im{\mathop{\mathrm{Im}}}
\def\Coker{\mathop{\mathrm{Coker}}}
\def\Coim{\mathop{\mathrm{Coim}}}

% Boxtimes
%
\DeclareMathSymbol{\boxtimes}{\mathbin}{AMSa}{"02}

%
% Macros for moduli stacks/spaces
%
\def\QCohstack{\mathcal{QC}\!\mathit{oh}}
\def\Cohstack{\mathcal{C}\!\mathit{oh}}
\def\Spacesstack{\mathcal{S}\!\mathit{paces}}
\def\Quotfunctor{\mathrm{Quot}}
\def\Hilbfunctor{\mathrm{Hilb}}
\def\Curvesstack{\mathcal{C}\!\mathit{urves}}
\def\Polarizedstack{\mathcal{P}\!\mathit{olarized}}
\def\Complexesstack{\mathcal{C}\!\mathit{omplexes}}
% \Pic is the operator that assigns to X its picard group, usage \Pic(X)
% \Picardstack_{X/B} denotes the Picard stack of X over B
% \Picardfunctor_{X/B} denotes the Picard functor of X over B
\def\Pic{\mathop{\mathrm{Pic}}\nolimits}
\def\Picardstack{\mathcal{P}\!\mathit{ic}}
\def\Picardfunctor{\mathrm{Pic}}
\def\Deformationcategory{\mathcal{D}\!\mathit{ef}}
\usepackage{cleveref}
\crefname{originalreference}{original reference}{original references}
\Crefname{originalreference}{Original reference}{Original references}
\makeatletter
\newlabel{decent-spaces-lemma-decent-generic-points}{{0ABV}{source}{}{sourcepointer-1}{}}
\newlabel{decent-spaces-lemma-decent-generic-points@cref}{{[originalreference][0][]0ABV}{[1][0][]source}{}{}{}}
\makeatother
\crefname{originalreference}{original reference}{original references}
\Crefname{originalreference}{Original reference}{Original references}
\makeatletter
\newlabel{decent-spaces-lemma-properties-trivial-implications}{{03M5}{source}{}{sourcepointer-2}{}}
\newlabel{decent-spaces-lemma-properties-trivial-implications@cref}{{[originalreference][0][]03M5}{[1][0][]source}{}{}{}}
\makeatother
\crefname{originalreference}{original reference}{original references}
\Crefname{originalreference}{Original reference}{Original references}
\makeatletter
\newlabel{decent-spaces-lemma-property-over-property}{{0ABY}{source}{}{sourcepointer-3}{}}
\newlabel{decent-spaces-lemma-property-over-property@cref}{{[originalreference][0][]0ABY}{[1][0][]source}{}{}{}}
\makeatother
\crefname{originalreference}{original reference}{original references}
\Crefname{originalreference}{Original reference}{Original references}
\makeatletter
\newlabel{morphisms-lemma-finite-birational-over-normal}{{0AB1}{source}{}{sourcepointer-4}{}}
\newlabel{morphisms-lemma-finite-birational-over-normal@cref}{{[originalreference][0][]0AB1}{[1][0][]source}{}{}{}}
\makeatother
\crefname{originalreference}{original reference}{original references}
\Crefname{originalreference}{Original reference}{Original references}
\makeatletter
\newlabel{spaces-cohomology-lemma-flat-base-change-cohomology}{{073K}{source}{}{sourcepointer-5}{}}
\newlabel{spaces-cohomology-lemma-flat-base-change-cohomology@cref}{{[originalreference][0][]073K}{[1][0][]source}{}{}{}}
\makeatother
\crefname{originalreference}{original reference}{original references}
\Crefname{originalreference}{Original reference}{Original references}
\makeatletter
\newlabel{spaces-morphisms-lemma-normalization-generic}{{0AYG}{source}{}{sourcepointer-6}{}}
\newlabel{spaces-morphisms-lemma-normalization-generic@cref}{{[originalreference][0][]0AYG}{[1][0][]source}{}{}{}}
\makeatother
\crefname{originalreference}{original reference}{original references}
\Crefname{originalreference}{Original reference}{Original references}
\makeatletter
\newlabel{spaces-morphisms-lemma-normalization-in-reduced}{{0AYF}{source}{}{sourcepointer-7}{}}
\newlabel{spaces-morphisms-lemma-normalization-in-reduced@cref}{{[originalreference][0][]0AYF}{[1][0][]source}{}{}{}}
\makeatother
\crefname{originalreference}{original reference}{original references}
\Crefname{originalreference}{Original reference}{Original references}
\makeatletter
\newlabel{spaces-morphisms-lemma-properties-normalization}{{0ABP}{source}{}{sourcepointer-8}{}}
\newlabel{spaces-morphisms-lemma-properties-normalization@cref}{{[originalreference][0][]0ABP}{[1][0][]source}{}{}{}}
\makeatother
\crefname{originalreference}{original reference}{original references}
\Crefname{originalreference}{Original reference}{Original references}
\makeatletter
\newlabel{spaces-over-fields-definition-integral-algebraic-space}{{0AD4}{source}{}{sourcepointer-9}{}}
\newlabel{spaces-over-fields-definition-integral-algebraic-space@cref}{{[originalreference][0][]0AD4}{[1][0][]source}{}{}{}}
\makeatother
\crefname{originalreference}{original reference}{original references}
\Crefname{originalreference}{Original reference}{Original references}
\makeatletter
\newlabel{spaces-over-fields-lemma-normal-integral-cover-by-affines}{{0AYH}{source}{}{sourcepointer-10}{}}
\newlabel{spaces-over-fields-lemma-normal-integral-cover-by-affines@cref}{{[originalreference][0][]0AYH}{[1][0][]source}{}{}{}}
\makeatother
\crefname{originaltheorem}{original theorem}{original theorems}
\Crefname{originaltheorem}{Original theorem}{Original theorems}
\makeatletter
\newlabel{theorem-stein-factorization-general}{{0A1C}{source}{}{sourcepointer-11}{}}
\newlabel{theorem-stein-factorization-general@cref}{{[originaltheorem][0][]0A1C}{[1][0][]source}{}{}{}}
\makeatother
\title{Geometrically connected fibres of a proper morphism towards a normal integral base}
\author{The Stacks Project Authors}
\date{stacks-project:spaces-more-morphisms.tex:local-snapshot-20261011; GNU FDL 1.2 or later}
\begin{document}
\maketitle
\noindent\textit{Editorial note.} Counted span is the author's own proof environment, source lines 8304-8344, immediately following the statement at 8286-8302.\par\medskip
\noindent\textit{Editorial note.} Same-name pair note: this is the spaces (more general) side of a label that also exists in chapter more-morphisms; the schemes side is not counted (see dedup.scope).\par\medskip
\noindent\textit{Editorial note.} Three cross-chapter pointer families (Decent Spaces, Cohomology of Spaces, Spaces over Fields) are recorded in assembly.external\_references; all same-chapter pointers are line-indexed stubs into the pinned source.\par\medskip
\noindent\textit{Allowed standard prerequisites (editorial).}\begin{itemize}
\item Stein factorization for algebraic spaces. Theorem theorem-stein-factorization-general of this chapter, cited at proof line 8304. Provides the factorization X -> Y' -> Y with Y' the relative spectrum of f\_*O\_X; the whole proof is the identification Y' = Y.
\item Etale-local nature of the normalization and of Y'. Morphisms of Spaces, Lemma spaces-morphisms-lemma-properties-normalization, cited at proof lines 8310-8312. Formation of Y' commutes with etale base change, which lets the problem be moved to an affine integral etale cover V -> Y.
\item Covering a normal integral algebraic space by affines. Spaces over Fields, Lemma spaces-over-fields-lemma-normal-integral-cover-by-affines, cited at proof lines 8308-8310. Reduces the isomorphism Y' = Y to checking Y' x\_Y V -> V over affine integral schemes V.
\item Finite birational morphisms towards a normal scheme are isomorphisms. Morphisms, Lemma morphisms-lemma-finite-birational-over-normal, cited at proof lines 8326-8328. The criterion applied to Y' -> Y once both are normal integral schemes with the same function field.
\item Decentness of X over a decent separated base. Decent Spaces, Lemmas decent-spaces-lemma-properties-trivial-implications and decent-spaces-lemma-property-over-property, cited at proof lines 8332-8335. Transfers decency from Y to X so that the generic-point bookkeeping of assumptions (5) and (6) applies to |X|.
\item Flat base change for cohomology. Cohomology of Spaces, Lemma spaces-cohomology-lemma-flat-base-change-cohomology, cited at proof lines 8319-8321. Shows condition (6) (H\textasciicircum{}0 of the generic fibre equals kappa(xi)) is preserved by the etale base change V -> Y.
\end{itemize}\medskip
\ifdefined\hypertarget\hypertarget{sourcepointer-1}{}\fi
\noindent\textit{Source reference.} Original reference 0ABV: focused-seeds/geometry/S9-spaces-geom-conn-fibres-towards-normal/ref-tags.txt, source lines 5--5. Navigation pointer only: the counted proof cites this statement by \textbackslash{}ref as a standing input; the deduction the counted proof has to make is carried out inside the counted span, and no part of that deduction is delegated to this pointer.\par\medskip
\ifdefined\hypertarget\hypertarget{sourcepointer-2}{}\fi
\noindent\textit{Source reference.} Original reference 03M5: focused-seeds/geometry/S9-spaces-geom-conn-fibres-towards-normal/ref-tags.txt, source lines 6--6. Navigation pointer only: the counted proof cites this statement by \textbackslash{}ref as a standing input; the deduction the counted proof has to make is carried out inside the counted span, and no part of that deduction is delegated to this pointer.\par\medskip
\ifdefined\hypertarget\hypertarget{sourcepointer-3}{}\fi
\noindent\textit{Source reference.} Original reference 0ABY: focused-seeds/geometry/S9-spaces-geom-conn-fibres-towards-normal/ref-tags.txt, source lines 7--7. Navigation pointer only: the counted proof cites this statement by \textbackslash{}ref as a standing input; the deduction the counted proof has to make is carried out inside the counted span, and no part of that deduction is delegated to this pointer.\par\medskip
\ifdefined\hypertarget\hypertarget{sourcepointer-4}{}\fi
\noindent\textit{Source reference.} Original reference 0AB1: focused-seeds/geometry/S9-spaces-geom-conn-fibres-towards-normal/ref-tags.txt, source lines 8--8. Navigation pointer only: the counted proof cites this statement by \textbackslash{}ref as a standing input; the deduction the counted proof has to make is carried out inside the counted span, and no part of that deduction is delegated to this pointer.\par\medskip
\ifdefined\hypertarget\hypertarget{sourcepointer-5}{}\fi
\noindent\textit{Source reference.} Original reference 073K: focused-seeds/geometry/S9-spaces-geom-conn-fibres-towards-normal/ref-tags.txt, source lines 9--9. Navigation pointer only: the counted proof cites this statement by \textbackslash{}ref as a standing input; the deduction the counted proof has to make is carried out inside the counted span, and no part of that deduction is delegated to this pointer.\par\medskip
\ifdefined\hypertarget\hypertarget{sourcepointer-6}{}\fi
\noindent\textit{Source reference.} Original reference 0AYG: focused-seeds/geometry/S9-spaces-geom-conn-fibres-towards-normal/ref-tags.txt, source lines 10--10. Navigation pointer only: the counted proof cites this statement by \textbackslash{}ref as a standing input; the deduction the counted proof has to make is carried out inside the counted span, and no part of that deduction is delegated to this pointer.\par\medskip
\ifdefined\hypertarget\hypertarget{sourcepointer-7}{}\fi
\noindent\textit{Source reference.} Original reference 0AYF: focused-seeds/geometry/S9-spaces-geom-conn-fibres-towards-normal/ref-tags.txt, source lines 11--11. Navigation pointer only: the counted proof cites this statement by \textbackslash{}ref as a standing input; the deduction the counted proof has to make is carried out inside the counted span, and no part of that deduction is delegated to this pointer.\par\medskip
\ifdefined\hypertarget\hypertarget{sourcepointer-8}{}\fi
\noindent\textit{Source reference.} Original reference 0ABP: focused-seeds/geometry/S9-spaces-geom-conn-fibres-towards-normal/ref-tags.txt, source lines 12--12. Navigation pointer only: the counted proof cites this statement by \textbackslash{}ref as a standing input; the deduction the counted proof has to make is carried out inside the counted span, and no part of that deduction is delegated to this pointer.\par\medskip
\ifdefined\hypertarget\hypertarget{sourcepointer-9}{}\fi
\noindent\textit{Source reference.} Original reference 0AD4: focused-seeds/geometry/S9-spaces-geom-conn-fibres-towards-normal/ref-tags.txt, source lines 13--13. Navigation pointer only: the counted proof cites this statement by \textbackslash{}ref as a standing input; the deduction the counted proof has to make is carried out inside the counted span, and no part of that deduction is delegated to this pointer.\par\medskip
\ifdefined\hypertarget\hypertarget{sourcepointer-10}{}\fi
\noindent\textit{Source reference.} Original reference 0AYH: focused-seeds/geometry/S9-spaces-geom-conn-fibres-towards-normal/ref-tags.txt, source lines 14--14. Navigation pointer only: the counted proof cites this statement by \textbackslash{}ref as a standing input; the deduction the counted proof has to make is carried out inside the counted span, and no part of that deduction is delegated to this pointer.\par\medskip
\ifdefined\hypertarget\hypertarget{sourcepointer-11}{}\fi
\noindent\textit{Source reference.} Original theorem 0A1C: sources/stacks-spaces-more-morphisms/pinned-original-source.tex, source lines 8197--8197. Navigation pointer only: the counted proof cites this statement by \textbackslash{}ref as a standing input; the deduction the counted proof has to make is carried out inside the counted span, and no part of that deduction is delegated to this pointer.\par\medskip
\begin{lemma}
\label{lemma-geometrically-connected-fibres-towards-normal}
Let $S$ be a scheme. Let $f : X \to Y$ be a morphism of algebraic
spaces over $S$. Assume
\begin{enumerate}
\item $f$ is proper,
\item $Y$ is integral (Spaces over Fields, Definition
\ref{spaces-over-fields-definition-integral-algebraic-space})
with generic point $\xi$,
\item $Y$ is normal,
\item $X$ is reduced,
\item every generic point of an irreducible component of $|X|$ maps to $\xi$,
\item we have $H^0(X_\xi, \mathcal{O}) = \kappa(\xi)$.
\end{enumerate}
Then $f_*\mathcal{O}_X = \mathcal{O}_Y$ and $f$
has geometrically connected fibres.
\end{lemma}

\begin{proof}
Apply Theorem \ref{theorem-stein-factorization-general} to get a
factorization $X \to Y' \to Y$. It is enough to show that $Y' = Y$.
It suffices to show that $Y' \times_Y V \to V$ is an isomorphism,
where $V \to Y$ is an \'etale morphism and $V$ an affine integral scheme,
see Spaces over Fields, Lemma
\ref{spaces-over-fields-lemma-normal-integral-cover-by-affines}.
The formation of $Y'$ commutes with \'etale base change, see
Morphisms of Spaces, Lemma
\ref{spaces-morphisms-lemma-properties-normalization}.
The generic points of $X \times_Y V$ lie over the generic points of $X$
(Decent Spaces, Lemma \ref{decent-spaces-lemma-decent-generic-points})
hence map to the generic point of $V$ by assumption (5). Moreover, condition
(6) is preserved under the base change by $V \to Y$, for example by
flat base change (Cohomology of Spaces, Lemma
\ref{spaces-cohomology-lemma-flat-base-change-cohomology}).
Thus it suffices to prove the lemma in case $Y$ is a normal
integral affine scheme.

\medskip\noindent
Assume $Y$ is a normal integral affine scheme. We will show $Y' \to Y$
is an isomorphism by an application of Morphisms, Lemma
\ref{morphisms-lemma-finite-birational-over-normal}.
Namely, $Y'$ is reduced because $X$ is reduced (Morphisms of Spaces, Lemma
\ref{spaces-morphisms-lemma-normalization-in-reduced}).
The morphism $Y' \to Y$ is integral by the theorem cited above.
Since $Y$ is decent and $X \to Y$ is separated, we see that
$X$ is decent too; to see this use
Decent Spaces, Lemmas
\ref{decent-spaces-lemma-properties-trivial-implications} and
\ref{decent-spaces-lemma-property-over-property}. By assumption (5),
Morphisms of Spaces, Lemma \ref{spaces-morphisms-lemma-normalization-generic},
and Decent Spaces, Lemma \ref{decent-spaces-lemma-decent-generic-points}
we see that every generic point of an irreducible component of $|Y'|$
maps to $\xi$. On the other hand, since $Y'$ is the relative
spectrum of $f_*\mathcal{O}_X$ we see that the scheme theoretic fibre
$Y'_\xi$ is the spectrum of $H^0(X_\xi, \mathcal{O})$ which is
equal to $\kappa(\xi)$ by assumption. Hence $Y'$ is an integral
scheme with function field equal to the function field of $Y$.
This finishes the proof.
\end{proof}

\end{document}
